% !TEX program = pdflatex
% ================================================================
% Writing sample (single-author technical note)
% Topic: COPF (deployment-stable counterfactual fairness) + heavy tails
% ================================================================

\documentclass[11pt]{article}

\usepackage[margin=1in]{geometry}
\usepackage{amsmath,amssymb,amsthm,mathtools}
\usepackage{bm}
\usepackage{microtype}
\usepackage{enumitem}
\usepackage{booktabs}
\usepackage{algorithm}
\usepackage{algpseudocode}
\usepackage{hyperref}
\usepackage{cleveref}

\hypersetup{colorlinks=true, linkcolor=blue, citecolor=blue, urlcolor=blue}

% --- notation ---
\newcommand{\E}{\mathbb{E}}
\newcommand{\Pp}{\mathbb{P}}
\newcommand{\Var}{\mathrm{Var}}
\newcommand{\Cov}{\mathrm{Cov}}
\newcommand{\1}{\mathbf{1}}
\newcommand{\R}{\mathbb{R}}
\newcommand{\cA}{\mathcal{A}}
\newcommand{\cH}{\mathcal{H}}
\newcommand{\cW}{\mathcal{W}}
\newcommand{\cF}{\mathcal{F}}
\newcommand{\cI}{\mathcal{I}}
\newcommand{\abs}[1]{\left|#1\right|}
\newcommand{\norm}[1]{\left\|#1\right\|}
\newcommand{\inner}[2]{\langle #1,#2\rangle}

\newcommand{\Eb}{\widehat{\mathbb{E}}} % self-normalized empirical average
\newcommand{\tGADR}{\widetilde{\Gamma}}

% --- cleveref names ---
\crefname{theorem}{Theorem}{Theorems}
\Crefname{theorem}{Theorem}{Theorems}
\crefname{lemma}{Lemma}{Lemmas}
\Crefname{lemma}{Lemma}{Lemmas}
\crefname{assumption}{Assumption}{Assumptions}
\Crefname{assumption}{Assumption}{Assumptions}
\crefname{definition}{Definition}{Definitions}
\Crefname{definition}{Definition}{Definitions}
\crefname{algorithm}{Algorithm}{Algorithms}
\Crefname{algorithm}{Algorithm}{Algorithms}

% --- theorem envs ---
\theoremstyle{plain}
\newtheorem{theorem}{Theorem}[section]
\newtheorem{lemma}[theorem]{Lemma}
\newtheorem{proposition}[theorem]{Proposition}
\newtheorem{corollary}[theorem]{Corollary}

\theoremstyle{definition}
\newtheorem{definition}[theorem]{Definition}
\newtheorem{assumption}[theorem]{Assumption}

\theoremstyle{remark}
\newtheorem{remark}[theorem]{Remark}

\newenvironment{proofsketch}{\begin{proof}[Proof sketch]}{\end{proof}}

% ================================================================
\title{\vspace{-1.2em}Writing Sample: Robust Certificates for Deployment-Stable\\ Counterfactual Fairness in Evolving Graph Recommendations}
\author{\vspace{-0.6em}Single-author technical note}
\date{\vspace{-0.8em}\today}

\begin{document}
\maketitle

\begin{abstract}
Online link recommendation on evolving graphs is \emph{performative}: changing which candidate links are shown
changes which links form and what feedback is observed later.
COPF (Counterfactual Online Performative Fairness) is a decision-layer framework that makes fairness
\emph{deployment-stable} by defining fairness on \emph{exposure counterfactuals} (shown vs.~not shown)
and certifying group gaps via Residual Outcome Indistinguishability (Residual-OI) on
graph-aware doubly-robust (GA-DR) pseudo-outcomes.

This note isolates a theory module I developed to strengthen COPF for two ``real-world'' failure modes:
(1) \textbf{heavy-tailed pseudo-outcomes} caused by inverse propensity weighting (small exposure probabilities),
and (2) \textbf{heavy-tailed graph degrees} (power-law networks), which invalidate assumptions that depend on the
maximum graph degree.

I show how to build \emph{robust} Residual-OI certificates under only a finite $p$-th moment assumption ($1<p\le 2$),
with explicit dependence on a \emph{checkable} interference degree $\kappa$ and mixing time $\tau_{\mathrm{mix}}$.
I then prove that in COPF, $\kappa$ can be upper-bounded \emph{by protocol design} (candidate budget and memory horizon),
so the guarantees remain meaningful even on hub-dominated (power-law) graphs.
\end{abstract}

\paragraph{What this writing sample is (and is not).}
This is a single-author technical note intended to be readable without the full COPF paper.
It summarizes the core COPF objects only to the extent needed for the two results above.

\paragraph{Main takeaways.}
\begin{itemize}[leftmargin=1.4em]
  \item \textbf{Robust auditing/certificates:} replace ``bounded residual'' assumptions (or heuristic clipping) with
  a principled robust mean estimator that tolerates heavy tails and dependence.
  \item \textbf{Hubs are not the bottleneck:} the statistical dependence parameter entering the certificate is not
  the maximum graph degree; it is the maximum degree $\kappa$ of a \emph{dependency graph over logged instances},
  which can be bounded by the logging protocol.
\end{itemize}

% ================================================================
\section{Problem: performative recommendation and deployment-stable fairness}
\label{sec:intro}
% ================================================================

In recommendation-driven graph growth (``who to follow'', ``people you may know''),
the platform selects a slate of candidate destinations each round.
Many edges form only after exposure, so the platform's policy affects the future graph and the future data.
As a result, fairness metrics computed directly on realized outcomes can drift after a policy update.

COPF addresses this by defining fairness on \emph{exposure counterfactuals}:
for each candidate dyad, compare what would happen if it were shown versus not shown,
and compare these effects across protected groups.
To estimate these counterfactual quantities from a single online stream,
COPF uses explicit exploration and propensity logging, and then audits residuals with
Outcome Indistinguishability-style tools.

\paragraph{Why heavy tails matter in COPF.}
Even if outcomes are binary, COPF uses inverse propensity weighting.
When some candidates have very small exposure probability $e_t(v)$,
the weight $1/e_t(v)$ can induce heavy-tailed pseudo-outcomes and residuals.
Separately, evolving social graphs often have power-law degrees;
any analysis that depends on the \emph{maximum} graph degree is therefore brittle.

This note develops a robust certificate that depends instead on (i) a finite moment assumption and
(ii) a dependence parameter that is controlled by the protocol.

% ================================================================
\section{Setup and COPF recap (minimal)}
\label{sec:setup}
% ================================================================

\paragraph{Rounds and candidates.}
At each round $t\in\{1,\dots,T\}$, the stream reveals a source node $u_t$ and a realized destination $v_t^\star$.
A candidate set is formed online as
\[C_t = \{v_t^\star\}\cup N_t,\]
where $N_t$ contains up to $N$ negatives (sampled without leakage).
Let $M:=\max_t |C_t|$ denote the (protocol) candidate budget.

\paragraph{Scores and exposure.}
A backbone scorer produces a probability $\hat p_t(u_t,v)\in(0,1)$ for each $v\in C_t$.
An exposure policy selects a size-$K$ slate $D_t\subseteq C_t$.
Let $D_t(v):=\1\{v\in D_t\}$ be the marginal exposure.
COPF logs each candidate's marginal entry propensity
\[e_t(v) := \Pp(D_t(v)=1\mid H_t),\]
where $H_t$ is the information used to construct the slate (graph history, scores, etc.).

\paragraph{Potential outcomes and counterfactual effects.}
For each candidate dyad $(u_t,v)$, let $Y_t^{(1)}(u_t,v)$ and $Y_t^{(0)}(u_t,v)$ denote
potential outcomes under exposure and no exposure.
The observed outcome satisfies consistency:
\[Y_t(u_t,v) = D_t(v)\,Y_t^{(1)}(u_t,v) + (1-D_t(v))\,Y_t^{(0)}(u_t,v).\]

\paragraph{Groups and slices.}
Each node $i$ has a protected attribute $A_i\in\cA$ (finite).
A dyad label $A_t(u_t,v)$ is defined by an attribution rule (source-side $A_{u_t}$ or destination-side $A_v$).
For auditing, we use a finite family $\cH$ of slice indicators
(e.g., group $\times$ score bucket $\times$ optional structural role).

\begin{assumption}[Identifiability and bounded dependence]\label{ass:ident}
For each round $t$ and candidate $v\in C_t$, define a logged context $W_t=W_t(u_t,v)$ computed only from $G_{\le t}$.
\begin{enumerate}[leftmargin=2em]
\item \textbf{Overlap:} $\exists\,e_{\min}>0$ such that $e_{\min}\le e_t(v)\le 1-e_{\min}$.
\item \textbf{Local ignorability:} conditional on $W_t$, $D_t(v)$ is independent of $(Y_t^{(0)}(u_t,v),Y_t^{(1)}(u_t,v))$.
\item \textbf{Bounded dependence:} the process exhibits temporal dependence with mixing time $\tau_{\mathrm{mix}}$ and
local network interference that induces a dependency graph of maximum degree $\kappa$ on any evaluation window.
\end{enumerate}
\end{assumption}

\paragraph{GA-DR pseudo-outcomes.}
Let $\mu_a(W):=\E[Y^{(a)}\mid W]$ and let $\hat\mu_{a,t}(W)$ be online estimates.
For $a\in\{0,1\}$ define
\begin{equation}
\label{eq:gadr}
\tGADR_t^{(a)}(u_t,v)
:=\hat\mu_{a,t}(W_t)
+\frac{\1\{D_t(v)=a\}}{\hat e_t^{(a)}(v)}\bigl(Y_t(u_t,v)-\hat\mu_{a,t}(W_t)\bigr),
\end{equation}
where $\hat e_t^{(1)}(v)=\hat e_t(v)$ and $\hat e_t^{(0)}(v)=1-\hat e_t(v)$.

\paragraph{Residuals.}
For calibration we use
\[\hat r_t^{(0)}(u_t,v):=\tGADR_t^{(0)}(u_t,v)-\hat p_t(u_t,v),\]
and for treatment-effect parity we use the effect residual
\[\hat r_t^{(\Delta)}(u_t,v):=\bigl(\tGADR_t^{(1)}(u_t,v)-\tGADR_t^{(0)}(u_t,v)\bigr)-\hat\tau_t(W_t),\]
where $\hat\tau_t(W)$ is an optional centering term.

\paragraph{Graph-aware window average.}
Let $w_t(u_t,v)\ge 0$ be (optional) graph-aware weights.
On a window $\cW$ of logged instances, define the self-normalized mean
\[
\Eb_{\mathrm{GA},\cW}[Z] :=
\frac{\sum_{(t,v)\in\cW} w_t(u_t,v)\,Z_t(u_t,v)}{\sum_{(t,v)\in\cW} w_t(u_t,v)}.
\]

% ================================================================
\section{Residual-OI and certificate reduction}
\label{sec:cert}
% ================================================================

\begin{definition}[Windowed Residual-OI]\label{def:oi}
Given residuals $\hat r_t$ and an auditor family $\cH$ of bounded tests $h_t=h(W_t)\in[-1,1]$,
\[
\widehat{\mathrm{OI}}_{\cW}(\hat r;\cH) := \sup_{h\in\cH}\abs{\Eb_{\mathrm{GA},\cW}[h_t\hat r_t]}.
\]
\end{definition}

\paragraph{Fairness gaps as residual correlations.}
A key simplification in COPF is that counterfactual fairness gaps can be written as
(normalized) residual correlations on slices.
For example, define the exposure ATE for group $s$ as
$\tau_s := \E[Y^{(1)}-Y^{(0)}\mid A=s]$ and define the parity gap
$g_{\mathrm{TEgap}}:=\max_{s,s'}\abs{\tau_s-\tau_{s'}}$.
Similarly define a slice calibration gap for bucket $I$.

\begin{lemma}[Gap linearization to residual moments]\label{lem:linearize}
Assume overlap and $\Pp(A=s)>0$.
Let $r^{(\Delta)}:=(Y^{(1)}-Y^{(0)})-\tau(W)$ and $r^{(0)}:=Y^{(0)}-\hat p$.
Then for any groups $s_1,s_2$,
\[
\tau_{s_1}-\tau_{s_2}
= \frac{\E[\1\{A=s_1\}r^{(\Delta)}]}{\Pp(A=s_1)}
- \frac{\E[\1\{A=s_2\}r^{(\Delta)}]}{\Pp(A=s_2)}
+ \Bigl(\E[\tau(W)\mid A=s_1]-\E[\tau(W)\mid A=s_2]\Bigr),
\]
and for any calibration slice $(s,I)$,
\[
\E[Y^{(0)}\mid A=s,\hat p\in I]-\E[\hat p\mid A=s,\hat p\in I]
= \frac{\E[\1\{A=s,\hat p\in I\}r^{(0)}]}{\Pp(A=s,\hat p\in I)}.
\]
\end{lemma}

\begin{remark}[Certificate form]
If $\tau(W)$ is group-agnostic (as in many COPF instantiations),
the last term in \Cref{lem:linearize} cancels between groups.
Then bounding $\sup_{h\in\cH}\abs{\E[h r]}$ and dividing by a minimum slice mass $p_{\min}$
yields a uniform certificate for both parity and calibration gaps.
\end{remark}

So far, nothing addresses heavy tails:
standard concentration arguments typically assume $\hat r_t\in[-1,1]$ (via clipping).
Next, we remove this boundedness requirement.

% ================================================================
\section{Robust Residual-OI under heavy tails and dependence}
\label{sec:heavy}
% ================================================================

\subsection{Why heavy tails occur}
Even with $Y\in\{0,1\}$, the GA-DR term \eqref{eq:gadr} contains $1/\hat e_t^{(a)}(v)$.
When some $\hat e_t(v)$ are small, $\tGADR$ and hence residuals can have long tails.
Clipping $1/\hat e$ fixes this numerically but complicates theory.
Instead we assume a finite moment and use a robust mean estimator.

\subsection{Assumptions: $p$-moment heavy tails}
\begin{assumption}[$p$-moment residuals]\label{ass:pmoment}
Fix $p\in(1,2]$.
For each auditor $h\in\cH$, define the scalar process
$X_{t,h}:=h_t\,r_t$ where $r_t$ is either $r_t^{(0)}$ or $r_t^{(\Delta)}$.
Assume a uniform conditional $p$-moment bound
\[
\E\bigl[\abs{X_{t,h}}^p\mid \cF_{t-1}\bigr] \le \sigma^p
\quad\text{for all }t\text{ and all }h\in\cH.
\]
Additionally, assume the bounded dependence parameters $(\kappa,\tau_{\mathrm{mix}})$ from \Cref{ass:ident}.
\end{assumption}

\subsection{Estimator: truncated self-normalized correlations}
Fix a window $\cW$ with $|\cW|=T$ and a truncation level $\tau>0$.
Define the truncated window statistic
\[
\widehat m_{\tau}(h)
:= \Eb_{\mathrm{GA},\cW}\bigl[X_{t,h}\,\1\{\abs{X_{t,h}}\le \tau\}\bigr].
\]
This is an estimator of $\E[X_{t,h}]$ with an explicit bias--variance tradeoff.

\subsection{Main theorem: uniform concentration with dependence corrections}
\begin{theorem}[Robust auditing bound for heavy-tailed COPF residuals]\label{thm:robust}
Assume \Cref{ass:pmoment} and let $\cH$ be finite.
Let $T_{\mathrm{eff}}$ denote an effective sample size satisfying
\begin{equation}
\label{eq:teff}
T_{\mathrm{eff}} \asymp \frac{T}{(\kappa+1)\,\tau_{\mathrm{mix}}}
\qquad\text{(up to universal constants and logarithmic factors).}
\end{equation}
Choose
\begin{equation}
\label{eq:tauchoice}
\tau \asymp \Bigl(\frac{\sigma^p\,T_{\mathrm{eff}}}{\log(\abs{\cH}/\delta)}\Bigr)^{1/p}
\end{equation}
for a confidence level $\delta\in(0,1)$.
Then with probability at least $1-\delta$,
\[
\sup_{h\in\cH}\abs{\widehat m_{\tau}(h)-\E[X_{t,h}]}
\ \le\ C\,\sigma\Bigl(\frac{\log(\abs{\cH}/\delta)}{T_{\mathrm{eff}}}\Bigr)^{\frac{p-1}{p}},
\]
for a universal constant $C>0$.

Consequently, if all audited slices have mass at least $p_{\min}$,
then the corresponding counterfactual fairness gaps that admit the ratio form in \Cref{lem:linearize}
obey a high-probability certificate
\[
\max\{g_{\mathrm{Calmax}},\ g_{\mathrm{TEgap}}\}
\ \lesssim\ \frac{1}{p_{\min}}\,\sigma\Bigl(\frac{\log(\abs{\cH}/\delta)}{T_{\mathrm{eff}}}\Bigr)^{\frac{p-1}{p}}
\quad\text{(up to a factor of $2$ for two-group differences).}
\]
\end{theorem}

\begin{proofsketch}[Proof sketch of \Cref{thm:robust}]
\textbf{(1) Truncation bias.}
By Markov and the $p$-moment bound,
$\abs{\E[X\1\{\abs{X}>\tau\}]}\le \E[\abs{X}^p]\,\tau^{-(p-1)}\le \sigma^p\tau^{-(p-1)}$.

\textbf{(2) Concentration of truncated terms.}
The truncated variables are bounded by $\tau$.
Within a window, bounded local interference yields a dependency graph of degree $\kappa$,
so a coloring argument reduces the sum to $\kappa+1$ nearly independent sub-sums.
Temporal mixing is handled by blocking at scale $\tau_{\mathrm{mix}}$.
This yields an effective number of approximately independent blocks $T_{\mathrm{eff}}$ as in \eqref{eq:teff}.
Applying a Hoeffding/Bernstein inequality for bounded (nearly) independent blocks and union bounding over $\cH$
controls the deviation of $\widehat m_{\tau}(h)$ from its truncated expectation.

\textbf{(3) Balance bias and deviation.}
Bounding the truncated second moment by $\E[(X\1\{\abs{X}\le \tau\})^2]\le \sigma^p\tau^{2-p}$
shows that the deviation term scales like
$\sigma^{p/2}\tau^{1-p/2}\sqrt{\log(\abs{\cH}/\delta)/T_{\mathrm{eff}}}$.
Choosing $\tau$ as in \eqref{eq:tauchoice} balances this deviation with the truncation bias,
yielding the final rate $T_{\mathrm{eff}}^{-(p-1)/p}$.
\end{proofsketch}

\begin{remark}[Interpretation]
When $p=2$ (finite variance), the rate becomes nearly $T_{\mathrm{eff}}^{-1/2}$ (up to logs),
matching sub-Gaussian behavior.
For heavier tails ($p<2$), the certificate degrades gracefully.
Importantly, dependence enters only through $T_{\mathrm{eff}}$, so controlling $\kappa$ is critical.
\end{remark}

% ================================================================
\section{Heavy-tailed graphs: protocol-level control of the interference degree $\kappa$}
\label{sec:kappa}
% ================================================================

This section addresses the ``hub'' issue.
In many social graphs, $\max_i \deg(i)$ is extremely large and can grow with $|V|$.
A naive analysis that upper bounds dependence by $\max_i\deg(i)$ becomes vacuous.

The key observation is that COPF's dependence parameter is not a graph degree.
It is the maximum degree of a \emph{dependency graph over logged instances} within a window.
This degree can be bounded by design choices in the logging protocol.

\subsection{Interference neighborhoods and dependency graphs}
Index a logged instance by $u=(t,v)$, i.e., a specific candidate $v\in C_t$.
Let $D$ denote the global exposure assignments over all instances.

\begin{definition}[Local interference neighborhood]\label{def:localint}
Fix integers $r_{\mathrm{int}}\ge 0$ and $L_{\mathrm{int}}\ge 0$.
We say the process satisfies $(r_{\mathrm{int}},L_{\mathrm{int}})$-local interference if
for each instance $u=(t,v)$ there exists a \emph{computable} set of instances $\cI(u)$ such that:
\begin{enumerate}[leftmargin=2em]
\item (Locality) $\cI(u)$ includes only instances from rounds $t'\in[t-L_{\mathrm{int}},\,t+L_{\mathrm{int}}]$;
\item (Graph locality) membership in $\cI(u)$ depends only on nodes within $r_{\mathrm{int}}$ hops of $\{u_t,v\}$ in $G_{\le t}$
(and the candidate rule);
\item (Interference restriction) the potential outcome of $u$ satisfies
$Y_u(D)=Y_u\bigl(g_u(D_{\cI(u)})\bigr)$ for some mapping $g_u$;
that is, treatments outside $\cI(u)$ do not affect $Y_u$.
\end{enumerate}
\end{definition}

\begin{definition}[Instance dependency graph and degree $\kappa$]\label{def:depgraph}
Fix a window $\cW$ of instances.
Define a graph on $\cW$ by connecting $u$ and $u'$ if $u'\in\cI(u)$ or $u\in\cI(u')$.
Let $\kappa$ denote the maximum degree of this dependency graph.
\end{definition}

\subsection{Main lemma: $\kappa$ is bounded by the candidate budget}
\begin{lemma}[Protocol-level bound on $\kappa$]\label{lem:kappa}
Assume $(r_{\mathrm{int}},L_{\mathrm{int}})$-local interference (\Cref{def:localint})
and a bounded candidate budget $|C_t|\le M$ for all rounds.
Then for any evaluation window, the induced dependency-graph degree satisfies
\[
\kappa \le (2L_{\mathrm{int}}+1)\,(M-1)\cdot c_0,
\]
where $c_0\ge 1$ is a small constant depending only on whether we index instances by candidates,
by dyads, or by per-round slates (i.e., bookkeeping conventions).
In particular, the bound is \emph{independent} of the maximum graph degree.
\end{lemma}

\begin{proofsketch}
Fix an instance $u=(t,v)$.
By definition, $\cI(u)$ can only include instances from rounds within $\pm L_{\mathrm{int}}$ of $t$.
Each such round contains at most $M$ instances (one per candidate).
Thus $\abs{\cI(u)}\le (2L_{\mathrm{int}}+1)\,M$.
The dependency graph degree of $u$ is at most $\abs{\cI(u)}-1$ up to constant-factor bookkeeping,
yielding the claim.
\end{proofsketch}

\begin{remark}[Why this resolves the ``power-law degree'' concern]
Even if the underlying evolving graph has hubs with degree $\Theta(|V|^\alpha)$,
the certificate in \Cref{thm:robust} remains meaningful as long as the protocol enforces
bounded candidate processing per round and bounded interference horizon.
This matches real systems, where candidate generation and logging are necessarily budgeted.
\end{remark}

% ================================================================
\section{How to use this in practice: a robust online Residual-OI module}
\label{sec:practice}
% ================================================================

\paragraph{Replace ``clip residuals'' by robust window statistics.}
In a COPF-style runner that audits a finite slice family $\cH$ on sliding windows,
one can implement \Cref{thm:robust} by computing truncated correlations.
A practical choice is to set $\tau$ using an empirical scale estimate per window
(e.g., MAD or a running $p$-moment proxy), while treating $p$ as a hyperparameter.

\begin{algorithm}[t]
\caption{Robust windowed Residual-OI (truncated correlations)}
\label{alg:robustoi}
\begin{algorithmic}[1]
\State \textbf{Input:} window $\cW$, auditor family $\cH$, truncation $\tau$, min-mass threshold $p_{\min}$
\For{each $h\in\cH$}
  \State compute $X_{t,h}=h_t\hat r_t$ for instances in $\cW$
  \State compute $\widehat m_{\tau}(h)=\Eb_{\mathrm{GA},\cW}[X_{t,h}\1\{\abs{X_{t,h}}\le \tau\}]$
\EndFor
\State output $\widehat{\mathrm{OI}}^{\mathrm{rob}}_{\cW}=\max_{h\in\cH}\abs{\widehat m_{\tau}(h)}$
\State (optional) report certificate $\widehat{\mathrm{OI}}^{\mathrm{rob}}_{\cW}/p_{\min}$
\end{algorithmic}
\end{algorithm}

\paragraph{Interpretation for a reader.}
Algorithm~\ref{alg:robustoi} is the exact point where ``heavy tails'' enter COPF.
The rest of the pipeline (identifiability, GA-DR, and gap linearization) remains unchanged.

% ================================================================
\section{Optional appendix: where the $p$-moment condition comes from}
\label{sec:appendix}
% ================================================================

A simple sufficient condition for \Cref{ass:pmoment} is that inverse propensity weights have bounded moments.
The next lemma connects this to exploration and the candidate budget.

\begin{lemma}[Inverse propensity moments under $\epsilon$-uniform Top-$K$ exploration]\label{lem:prop}
Suppose the exposure policy mixes in an $\epsilon$-uniform Top-$K$ component:
with probability $\epsilon$ it samples $K$ items uniformly without replacement from $C_t$.
Then for every $v\in C_t$,
\[
e_t(v)\ge \epsilon\,\frac{K}{|C_t|} \ge \epsilon\,\frac{K}{M}.
\]
Consequently, for any $p\ge 1$,
\[
\Bigl(\frac{1}{e_t(v)}\Bigr)^p \le \Bigl(\frac{M}{\epsilon K}\Bigr)^p.
\]
If outcomes are bounded (e.g., $Y\in[0,1]$) and $\hat\mu$ is bounded, then the GA-DR term
has finite $p$-moments with constants depending polynomially on $M/(\epsilon K)$.
\end{lemma}

\begin{remark}
This lemma makes a practical point: heavy tails are controlled primarily by (i) ensuring overlap through
explicit exploration and (ii) keeping candidate sets budgeted.
Neither requires assumptions on the maximum degree of the evolving graph.
\end{remark}

% ================================================================
\section*{References (minimal)}
\begin{thebibliography}{9}
\bibitem{copf}
COPF: An Online Framework for Deployment-Stable Counterfactual Fairness in Evolving Graphs.
(Under review; technical report version.)

\bibitem{oi}
Outcome indistinguishability and multicalibration references; see, e.g., Dwork et al.

\bibitem{dr}
Doubly robust estimation and off-policy evaluation references; see, e.g., Swaminathan \& Joachims.
\end{thebibliography}

\end{document}
